\allowdisplaybreaks
\begin{center}
	\hrule
	\vspace{.4cm}
 \begin{tabular*}{\textwidth}{@{}l@{}|@{\extracolsep{0.6in}}r@{}}%
\parbox{4.25in}{\raggedright{\includegraphics[width=.9\linewidth]{ictp-pwf.pdf}}} &
\parbox[c][]{4in}{{\Large\textbf{Md Akiful Islam Zawad} \par}
                    { Brac University \par}
                    { KHA 224, Progati Sarani, Merul Badda, \par}
                    { Dhaka 1212, Bangladesh \par}
                    { \href{ext.akiful.islam@bracu.ac.bd}{ext.akiful.islam@bracu.ac.bd}} \par}
\end{tabular*}\vspace{.3in}
\definecolor{ceruleanblue}{rgb}{0.16, 0.32, 0.6}
	\LARGE\scshape\textbf{\textcolor{ceruleanblue}{Physics for Bangladesh QFT School}}
\end{center}
\hrule\vspace{.25in}
{\large\textbf{Problem Sheet:}\ \textsc{3} \hspace{\hfill} \large\textbf{Due Date:} \today\\
	\hrule}
 %----------------------------
\paragraph*{Problem Set 1} %\hfill \newline
\\
The Lagrangian for the electromagnetic vector potential $A_\mu$, in the presence of a source current $j_\mu$, is given by
\begin{align}
    L = -\frac{1}{4} F_{\mu\nu} F^{\mu\nu} - j_\mu A^\mu, \label{eq:lagrangian-em-vector}
\end{align}
where $F_{\mu\nu} \equiv \partial_\mu A_\nu - \partial_\nu A_\mu$ is called the electromagnetic field strength tensor.
\begin{itemize}
    \item[(a)] Find the Euler-Lagrange equation corresponding to this Lagrangian.
    \item[(b)] The Euler-Lagrange equation will give rise to two of Maxwell’s equations. Recast these equations into Maxwell’s equations of electrodynamics.
\end{itemize}
\bigskip\bigskip\hline\hline\bigskip
\subsection*{Solution}
\begin{itemize}
\item [(a)] To find the Euler-Lagrange equation, we will use the following formula for a field theory:
\begin{align}
    \frac{\partial L}{\partial A^\mu} - \partial_\nu \left( \frac{\partial L}{\partial (\partial_\nu A^\mu)} \right) = 0.\label{eq:euler-lagrangian}
\end{align}
Compute the first term for the given Lagrangian (\ref{eq:lagrangian-em-vector}):

From the Lagrangian:
\begin{align*}
    L = -\frac{1}{4} F_{\mu\nu} F^{\mu\nu} - j_\mu A^\mu,
\end{align*}
the term involving $A^\mu$ explicitly is $-j_\mu A^\mu$. Thus,
\begin{align}
    \frac{\partial L}{\partial A^\mu} = -j_\mu.\label{eq:step-1}
\end{align}
Now compute the second term of (\ref{eq:euler-lagrangian}) for the given Lagrangian in (\ref{eq:lagrangian-em-vector}):

The derivative $\partial_\nu A_\mu$ only appears in $F_{\mu\nu}$. We need to express the Lagrangian in terms of $\partial_\nu A_\mu$. The Lagrangian contains the term:
\begin{align*}
    \frac{1}{4}F_{\mu\nu} F^{\mu\nu} = \frac{1}{4}(\partial_\mu A_\nu - \partial_\nu A_\mu)(\partial^\mu A^\nu - \partial^\nu A^\mu).
\end{align*}
Thus, $L$ depends on $\partial_\nu A_\mu$. To compute $\displaystyle\frac{\partial L}{\partial (\partial_\nu A_\mu)}$, let's focus on the first term:
\begin{align}
    -\frac{1}{4} F_{\mu\nu} F^{\mu\nu}= -\frac{1}{4} (\partial_\mu A_\nu - \partial_\nu A_\mu)(\partial^\mu A^\nu - \partial^\nu A^\mu).
\end{align}
We are interested in $\displaystyle\frac{\partial L}{\partial (\partial_\nu A_\mu)}$. Observe that:
\begin{align}
    \frac{\partial}{\partial (\partial_\nu A_\mu)} F_{\alpha\beta} = \delta^\nu_\alpha \delta^\mu_\beta - \delta^\nu_\beta \delta^\mu_\alpha.\label{eq:step-2}
\end{align}
Therefore:
\begin{align}
    \frac{\partial L}{\partial (\partial_\nu A^\mu)} = - F^{\nu\mu}.
\end{align}
Now we can plug (\ref{eq:step-1}) and (\ref{eq:step-2}) into (\ref{eq:euler-lagrangian}):
\begin{align}
    \frac{\partial L}{\partial A^\mu} - \partial_\nu \left( \frac{\partial L}{\partial (\partial_\nu A^\mu)} \right) &= 0\notag\\
    -j_\mu - \partial_\nu (-F^{\nu\mu}) &= 0,
\end{align}
which simplifies to:
\begin{align}
    \partial_\nu F^{\nu\mu} = j^\mu.
\end{align}
This is the Euler-Lagrange equation corresponding to the Lagrangian

\item [(b)] We will now show that this Euler-Lagrange equation corresponds to two of Maxwell’s equations.

Maxwell's equations in the presence of a current are:
\begin{align}
    \partial_\nu F^{\nu\mu} = j^\mu,
\end{align}
which we derived from the Euler-Lagrange equation.

Recall the components of $F_{\mu\nu}$:
\begin{align}
    F_{0i} &= E_i \quad &\text{(electric field components)}\label{eq:e-field-components}\\
    F_{ij} &= -\epsilon_{ijk} B_k\label{eq:b-field-components} \quad &\text{(magnetic field components)}
\end{align}
Let’s decompose this equation into its temporal and spatial components.

\textbf{Temporal component:} $\mu = 0$
\begin{align*}
    \partial_\nu F^{\nu 0} &= j^0\notag\\
    \partial_i F^{i0} &= j^0,
\end{align*}
or, using $F^{i0} = -E^i$:
\begin{align}
    \mathbf{\nabla} \cdot \mathbf{E} = j^0,
\end{align}
which is Gauss's law for Electric fields.

\textbf{Spatial component:} $\mu = i$
\begin{align*}
    \partial_\nu F^{\nu i} &= j^i\notag\\
    \partial_0 F^{0i} + \partial_j F^{ji} &= j^i.
\end{align*}
Using $F^{0i} = -E^i$ and $F^{ji} = \epsilon_{ijk} B^k$, this becomes:
\begin{align*}
    -\frac{\partial E^i}{\partial t} + \epsilon_{ijk} \partial_j B^k = j^i.
\end{align*}
This is the Ampère-Maxwell law:
\begin{align}
    \mathbf{\nabla} \times \mathbf{B} - \frac{\partial \mathbf{E}}{\partial t} = \mathbf{j}.
\end{align}
\end{itemize}
\bigskip\bigskip\hline\hline\bigskip